\chapter{Ra\'ul's decade}
\label{ch:raul}

\section{A different kind of governing}

Fidel Castro fell ill on 31 July 2006 and never resumed office; Ra\'ul Castro
took over provisionally that day and formally in February 2008. The change in
style was immediate and larger than the change in policy.

Fidel governed by speech: problems were named in public, solutions were
announced as campaigns, and the campaign was the instrument. Ra\'ul governed by
meeting. He delegated, set deadlines, demanded that ministries answer for
targets, and --- most unusually --- said in public that the system did not work.
The Camag\"uey speech of 26 July 2007 is the text: he told an audience that
Cuban wages were insufficient to live on, that structural and conceptual changes
were needed, and, in the example everyone remembers, that a country that could
not put a glass of milk on a child's table every day had a problem it should
stop explaining and start solving. Coming from the second man of the revolution,
in the anniversary speech, this was licence to discuss things that had not been
discussable.

He also opened a consultation. In late 2010 the draft economic guidelines were
circulated for discussion in workplaces and neighbourhoods; the government
reported some 160,000 meetings and millions of comments, and around two thirds
of the draft's paragraphs were modified \citep{mesalago2013}.\unv{} Whatever one makes of the
mechanism, this was the widest formal consultation the Cuban state had ever run,
and Cubans engaged with it seriously because they believed something would come
of it.

\section{The removal of humiliations}

The first measures were not economic reforms at all. They were the repeal of
prohibitions whose only function had been to remind Cubans of their position.

In 2008 Cubans were permitted to buy mobile telephones, computers, DVD players
and microwave ovens, and to stay in tourist hotels --- ending, twelve years
late, the situation in which a Cuban could not enter a building on Cuban soil
that a foreigner could. In 2011 they were permitted to buy and sell their own
houses and cars, which had been prohibited since 1960 and had produced an
elaborate grey market of swaps, fictional marriages and bribes.

And in October 2012, effective January 2013, the exit permit was abolished. The
\emph{tarjeta blanca} --- the white card, the permission-to-leave that had cost
money, taken months, required a letter of invitation, and been refused at
discretion --- was gone, and a Cuban could travel on a passport like anyone
else. Doctors and some other categories remained restricted in practice, and the
eight-year bar on returning was reduced to two. But the principle changed: after
fifty-four years, leaving Cuba temporarily stopped being a political act.

None of this cost the state anything. That is precisely the point, and
Chapter~\ref{ch:sequence} makes it an operating principle: the first phase of
any reform programme should begin with the prohibitions that cost nothing to
remove, because they are free, they are immediate, and they buy the credibility
that the expensive measures require. Ra\'ul Castro's most popular acts were all
of this kind.

\section{The guidelines, and why they stopped}

The Sixth Party Congress of April 2011 --- the first in fourteen years ---
adopted 313 Guidelines for Economic and Social Policy. The programme they
described was coherent: shrink the state payroll, expand non-state employment,
decentralise enterprise decisions, separate state administration from enterprise
management, unify the currency, and move from universal subsidy toward targeted
assistance. The announced first step was the transfer of half a million workers
off the state payroll within six months, and eventually more than a million.

The instruments were real. Self-employment was expanded to 178 and then 201
categories and the number of licence-holders rose from about 157,000 in 2010 to
over half a million by 2015 --- roughly a fifth of the employed
workforce \citep{ritter2015}.\unv{} Land was granted in usufruct under Decree-Law 259 of 2008,
initially 13 hectares, and under Decree-Law 300 of 2012 up to 67, to anyone
willing to farm it; something over 1.7 million hectares of idle land were handed
out.\unv{} Non-agricultural cooperatives were authorised on a pilot basis in
2013. The payroll reduction was never executed at anything like the announced
speed, because there was nowhere for the workers to go.

The reforms stalled, and the reason they stalled is the most useful thing in
this chapter for Part III, so it is worth stating precisely rather than
attributing to a general reluctance.

A Cuban private business between 2011 and 2021 had no legal personality. It was
a licensed individual, not a firm. It therefore could not sign a contract as an
entity, could not be sued or sue, could not own assets separately from its
owner, and could not, in most cases, be sold or inherited as a going concern.

It had no wholesale market. A restaurant bought its ingredients in the same
retail shops as households, at retail prices, competing with households for the
same scarce goods --- which is both economically absurd and the source of the
enduring popular complaint that the private sector caused the shortages.

It had no import rights. Everything had to be bought domestically or carried in
personal luggage by \emph{mulas}, travellers making repeated trips with
suitcases, which is how a substantial share of Cuban private-sector inputs
actually arrived for a decade.

It had no credit. Bank lending to the private sector was announced, hedged with
conditions, and never materialised at scale.

It had no access to foreign exchange at any usable rate.

And it operated by enumeration, not by right: an activity was legal if it
appeared on a list the state could shorten at any time, and the list was written
in the vocabulary of trades --- \emph{barber}, \emph{button-coverer},
\emph{party clown} --- with the professions excluded. A software developer had
no category. A design studio had no category.

The result is that Cuba got the private sector its rules were designed to
produce: several hundred thousand micro-operators in food, transport, lodging
and personal services, and essentially nothing in manufacturing, technology or
anything requiring an entity, a contract and a supply chain.
Chapter~\ref{ch:capital} treats this diagnosis as its starting point, because
almost every constraint in the list above is legal and internal, could be
removed without a dollar of foreign investment, and was not removed for a
decade.

\section{The opening}

On 17 December 2014 Ra\'ul Castro and Barack Obama announced simultaneously,
from Havana and Washington, that their countries would restore diplomatic
relations. The immediate exchange was the release of the American contractor
Alan Gross and of a Cuban-held intelligence asset, for the three remaining
members of the Cuban Five; behind it lay eighteen months of secret talks brokered
in Canada and by the Vatican \citep{leogrande2014}.

Embassies reopened in July 2015. Obama visited Havana in March 2016 --- the
first sitting American president since Coolidge in 1928 --- and gave a speech,
broadcast live on Cuban television, in which he said the embargo would end and
addressed Cubans directly about their own future. Regular flights resumed;
cruise ships docked; Airbnb entered the market in 2015 and had tens of thousands
of Cuban listings within two years; American visitors rose from around 91,000 in
2014 to over 600,000 in 2017, on top of a much larger flow of Cuban
Americans.\unv{}

For the Cuban private sector this was the best eighteen months it has ever had.
The demand went overwhelmingly to the non-state economy, because that is where
Americans on people-to-people licences were required to spend: private
restaurants, private rooms, private taxis, private guides. Money entered
households rather than the state, capital equipment came in with the visiting
relatives, and for the first time since 1959 there was a visible, growing,
domestically-owned Cuban middle class making its living from foreigners without
the state as intermediary.

That is also, precisely, why it was stopped.

\section{The brake}

The Seventh Party Congress of April 2016 was where the guidelines slowed. The
tone of the documents shifted from expansion of the non-state sector to its
regulation, and the phrase that circulated was the warning against the
concentration of wealth and property. In August 2017 the state suspended the
issue of new licences in a set of the most successful categories --- private
restaurants, room rentals, and others --- pending a review that lasted more than
a year and ended in a tighter regime: one licence per person, caps on the size
of restaurants, and stricter inspection.

Two further things happened in the same short period. Fidel Castro died on 25
November 2016. And on 8 November 2016 Donald Trump was elected, and in June 2017
issued a policy memorandum that reversed a substantial part of the opening ---
restricting individual people-to-people travel, prohibiting transactions with a
list of military-linked entities, and beginning the sequence of measures that
Chapter~\ref{ch:crisis} takes up.

It is tempting, and common, to attribute the closing entirely to Washington. The
sequence does not support it. The licence freeze of August 2017 was a Cuban
decision, taken about Cuban businesses, before the American measures had bitten.
Havana and Washington arrived at the same conclusion for opposite reasons:
that the growth of an independent Cuban private sector was politically
consequential. One of them wanted it and stopped subsidising it; the other did
not want it and stopped permitting it.

Ra\'ul Castro handed the presidency to Miguel D\'iaz-Canel in April 2018,
remaining First Secretary of the party until 2021. He left a country with a
larger private sector, a freer citizenry and a functioning relationship with the
United States, and with every structural problem identified in 2011 --- the dual
currency, the state payroll, enterprise autonomy, food production ---
diagnosed, announced, and unresolved.

\begin{ledger}{Ra\'ul's decade, 2006--2018}
\emph{Achieved.} The repeal of prohibitions that had no purpose but humiliation:
hotels, phones, computers, the sale of houses and cars, and above all the exit
permit, which after fifty-four years made leaving Cuba temporarily an ordinary
act. A private sector taken from 150,000 to over half a million people, about a
fifth of the workforce. Idle land distributed in usufruct on a large scale. A
national consultation the population took seriously. Diplomatic relations with
the United States restored, embassies reopened, and eighteen months in which
American demand flowed directly into Cuban households.

\emph{Cost.} A reform programme whose own instruments were designed so that it
could not succeed: no legal personality, no wholesale market, no import rights,
no credit, no foreign exchange, and legality by enumerated list. The result was
several hundred thousand micro-operators and nothing above them, which is
exactly what those rules produce. The one measure that would have made the
economy legible --- currency unification --- was announced in 2011 and deferred
for a decade, and would eventually be attempted at the worst possible moment.
And the opening that worked best was braked from the inside, in August 2017,
before Washington closed it from the outside.
\end{ledger}
